\documentclass[../elsevier.tex]{subfiles}
\begin{document}

\section{Related Work}
\label{sec:related}

\subsection{Inferring attributes from alters}
\label{sec:related-alters}

The operation we analyse has a long history in network analysis under several names.
In the network autocorrelation model an actor's outcome is regressed on a weighted
average of its alters' outcomes, and Leenders~\citeyearpar{leenders2002modeling} makes the
point on which this paper turns: the weight matrix is a substantive modelling
commitment, not a technical detail, and the conclusions depend on it. Row-normalised
weights---the natural choice, and the one implicit in most applications---are exactly
the mean aggregation of Section~\ref{sec:model}. Everett and
Borgatti~\citeyearpar{everett2026alter} give the general matrix formulation of \emph{alter
composition}, the profile of categories an actor is exposed to through its alters,
and extend it to overlapping memberships; our camouflage ratio is a two-category
alter-composition measure, and the question we ask is what happens to it when one
category chooses its ties.

Two in-journal results bear directly on the reliability of this inference.
Roeling and Nicholls~\citeyearpar{roeling2020imputation} impute unobserved attributes from
network structure and find the procedure sensitive to model misspecification and to
missingness. V\"or\"os et al.~\citeyearpar{voros2019limits} report a negative result of the
same kind: across five countries, friendship relations turn out to be a poor proxy
for status attributions, so one relation cannot be substituted for another. Both are
warnings that neighbour-based inference is conditional rather than automatic; our
contribution is to give one condition in closed form, for the case where the failure
is \emph{induced deliberately}.

Homophily is the assumption the inference rests on, and its measurement is a
developed literature here: Bojanowski and Corten~\citeyearpar{bojanowski2014measuring}
survey and axiomatise segregation and homophily indices, and
McMillan~\citeyearpar{mcmillan2022worth} shows that homophily measured on strong and weak
ties differs in magnitude, so that a single homophily figure for a population can be
misleading. Noel and Nyhan~\citeyearpar{noel2011unfriending} demonstrate the consequence for
inference, showing that homophily in tie \emph{retention} biases estimates of social
influence. Recent work in the journal continues in both directions: Sch\"u{\ss}ler and
Fuhse~\citeyearpar{schussler2025duality} argue that ties and attributes are mutually
constituting rather than one being read off the other, which is the conceptual
counterpart of our question; Espinosa-Rada et al.~\citeyearpar{espinosarada2026socioeconomic}
measure the socioeconomic composition of alters directly; and Tang et
al.~\citeyearpar{tang2026twomode} extend the autoregressive model to two-mode data with a
core--periphery degree structure, which is the setting our degree correction speaks to.
Where the network is only partially observed, Ajorlou et
al.~\citeyearpar{ajorlou2025dirichlet} give a sampling estimator for homophily itself.
Our results are the adversarial limit of this concern: we treat the degree
of homophily violation as a decision variable of one class rather than a property of
the population.

\subsection{Multiplexity, and why the analysis is per relation}
\label{sec:related-multiplex}

That different relations among the same actors have different structure is a founding
observation of this journal~\citep{lazega1999multiplexity}, and its methodological
consequence---that collapsing relations into a single adjacency destroys information,
while modelling them jointly quickly becomes intractable---is treated directly by
V\"or\"os and Snijders~\citeyearpar{voros2017cluster}. Snijders et
al.~\citeyearpar{snijders2013model} show the stronger version: structure apparently present
in an advice relation is substantially mediated by the friendship relation beneath
it. An et al.~\citeyearpar{an2026estimating} find the same phenomenon quantitatively, with
peer effects differing in magnitude between friends and classmates in a multilayer
setting.

Two recent contributions sharpen the point. Fluer et
al.~\citeyearpar{fluer2026survey} build multiplex models that carry interlayer
correlation explicitly rather than treating layers as independent, and Chandrasekhar et
al.~\citeyearpar{chandrasekhar2026multiplexing} show that multiplexity can have
opposite-signed consequences depending on the process running over it---inhibiting
low-reinforcement diffusion while enhancing high-reinforcement diffusion. That a
quantity can reverse sign across layer configurations is precisely what we find for the
usability of relational inference. Lakdawala et
al.~\citeyearpar{lakdawala2026bonds} meanwhile partition dyads into latent interaction
regimes, which is a complementary way of admitting that not all ties of a nominal type
are the same.

This matters for us empirically rather than only in principle.
Section~\ref{sec:measuring} finds camouflage differing by a factor approaching two
between relations recorded on the \emph{same} actors, which means our thresholds are
per-relation quantities and that a procedure pooling relations is averaging a usable
signal with an unusable one.

\subsection{Strategic tie formation}
\label{sec:related-strategic}

The premise that actors form ties instrumentally, and with an eye to how they will be
perceived, is established here but thinly populated. Burt et al.~\citeyearpar{burt2021network}
show experimentally that brokerage position causes attributions of leadership, and
report participants re-shaping an assigned network to suit their preferences. Shea et
al.~\citeyearpar{shea2019cheater} come closest to our mechanism: actors anticipating
unethical behaviour activate sparser networks that better conceal it, while those who
have already acted unethically activate denser ones for support---network structure
managed as a function of what the actor wants inferred about it. Barone and
Coscia~\citeyearpar{barone2018birds} estimate trustworthiness from mutually inconsistent
self-reports and find homophily in it, which is both an instance of inferring a
latent attribute that actors have an incentive to misreport and a demonstration that
the inference can work.

Nguyen~\citeyearpar{nguyen2026agency} adds the most recent instance, treating tie
formation as strategic behaviour conditional on structural position, and
Macdonald~\citeyearpar{macdonald2026wiretap} shows that covert networks carry
structural signatures distinct from ordinary collaboration---evidence that concealment
leaves traces in structure rather than erasing them.

What is absent, so far as we can determine, is the case where the tie portfolio is
\emph{engineered} to defeat the inference. We searched the journal's full record and
found no treatment of automated or inauthentic accounts; the nearest work concerns
misinformation susceptibility in personal networks~\citep{choi2026who}. We therefore
take our terminology from this journal---alter composition, autocorrelation,
exposure---and our substantive setting from the detection literature, rather than
importing the latter's vocabulary wholesale.

Within that detection literature a small strand does measure the structure of such
populations rather than only classifying them, and it is the strand closest to our
purpose: Ran et al.~\citeyearpar{ran2025identifying} characterise automated accounts by
motif frequencies across relation types, and Ariyarathne
et al.~\citeyearpar{ariyarathne2026behavior} treat measured behavioural change as the
discriminating signal in its own right. Two further results establish that camouflage is
an active and adaptive strategy rather than a fixed property: Pan et
al.~\citeyearpar{pan2026camera} document deliberate mimicry of benign users' content,
which breaks the assumption that structural and attribute cues agree, and Mukherjee
et al.~\citeyearpar{mukherjee2026optimal} show that detectors relying on aggregation can
be defeated by rewiring under realistic budget constraints. Our threshold gives the
condition under which such rewiring succeeds.

\subsection{Degree heterogeneity and the correction it forces}
\label{sec:related-degree}

Our edge-endpoint identity requires a degree correction, and the reason is a
long-standing concern here. Snijders~\citeyearpar{snijders1981degree} introduced degree
variance as an index of graph heterogeneity; Anderson et al.~\citeyearpar{anderson1999interaction}
showed that graph-level indices are systematically confounded by size and density; and
Krivitsky et al.~\citeyearpar{krivitsky2009representing} exhibit networks with identical,
highly skewed degree distributions but entirely different structure, concluding that
degree-based accounts cannot by themselves identify structure. Section~\ref{sec:measuring}
bears this out: the identity requires its degree factor wherever class mean degrees
differ, and holds exactly once it is included.

Kami\'nski et al.~\citeyearpar{kaminski2026relationships} extend the concern to
hypergraphs, finding an association between node degree and hyperedge size in all but a
small minority of $36$ empirical datasets, so degree is rarely separable from the
structure one wishes to measure.

The stochastic blockmodel we build on originates in this
journal~\citep{holland1983stochastic}. Degree correction \emph{within} blockmodels is
much less developed here; the one instance we found, Yu et
al.~\citeyearpar{yu2018detecting}, monitors node propensity in a dynamic degree-corrected
blockmodel in order to detect anomalous nodes, which is close to our setting in both
model and purpose.

\subsection{Analyses of neighbourhood averaging in the machine-learning literature}


The contextual stochastic block model~\citep{deshpande2018contextual} couples a
stochastic block model with Gaussian node attributes and has become the standard
setting for asking what message passing does to statistical separability. Baranwal
et al.~\citeyearpar{baranwal2021graph} established the central positive result: for a
two-class CSBM, one graph convolution improves the regime of linear separability by
a factor of roughly $1/\sqrt{D}$ in the expected degree $D$, and the resulting
classifier generalises out of distribution to graphs with different edge
probabilities.

That analysis assumes the generative process is homophilic---inter-class edges arise
from a fixed, non-adversarial probability---and establishes only the positive side.
Its follow-ups extend the positive results to multiple
layers~\citep{baranwal2023effects} and characterise the locally Bayes-optimal
architecture as one that \emph{interpolates} between an MLP in the low-graph-signal
regime and a convolution in the high-signal regime~\citep{baranwal2023optimality}.
Our Corollary~\ref{cor:critical} reduces to the separability gain
of~\citet{baranwal2021graph} when $\rho = 0$, which calibrates our model against this
baseline. That line has continued: Duranthon and
Zdeborov\'a~\citeyearpar{duranthon2025statistical} benchmark aggregation against the
information-theoretic optimum in the same generative model, quantifying how far short it
falls; Ma et al.~\citeyearpar{ma2025attention} show that attention---the closest
analogue to reweighting ties---helps only when structural noise dominates attribute
noise and hurts otherwise; and Baranwal~\citeyearpar{baranwal2026depth} obtains a sharp
dichotomy for the value of aggregation depth on sparse graphs, governed by a single
ratio. Our depth result in Section~\ref{sec:depth} is finite-$D$ and concerns the
crossover rather than the asymptotic regime, but it points the same way: depth is not
uniformly beneficial.

\paragraph{The crossover is known; our contribution is the imbalanced and
adversarial case.} We state the priority situation plainly, because it determines
what this paper does and does not claim. A closed-form threshold at which mean
aggregation becomes worse than the graph-free classifier was established by Ma
et al.~\citeyearpar[Thm.~2]{ma2022homophily} for the \emph{balanced} symmetric CSBM, and
Corollary~\ref{cor:critical} recovers their condition exactly at $c = 1$
(Remark~\ref{rem:ma} works the algebra through). Related crossover analyses have
since appeared: Luan et al.~\citeyearpar{luan2023when} compute the Bayes error of
graph-aware against graph-agnostic classifiers under a homophily parameter with
per-class variances and degrees, locating the intersection points directly and
criticising the fixed-classifier comparison used in~\citep{ma2022homophily}; Mao
et al.~\citeyearpar{mao2023demystifying} prove the per-node version, that a GCN beats an
MLP on homophilic nodes and loses on heterophilic ones \emph{within the same graph};
and the graph-attention analysis of Fountoulakis
et al.~\citeyearpar{fountoulakis2023graph} both identifies the same signal-deflation
factor and proves that attention---the closest analogue to soft edge
pruning---is lower-bounded by the better of convolution and the graph-free
classifier, while also showing (their Theorems~9 and~13) that no attention
mechanism separates intra- from inter-class edges in the hard regime.

Against this background our claim is narrower than ``we show when the graph is
harmful''. It is that the adversarial, class-imbalanced case is not a relabelling of
the balanced one: the two classes deflate asymmetrically, which moves the decision
boundary and invalidates a fixed-boundary comparison; bot prevalence enters the
threshold through $\eta = c\rho$ as a robustness parameter; and the mixture-variance
term omitted in~\citep{ma2022homophily} shifts the threshold by $11.5\%$. Read this way,
the closest antecedent to Theorem~\ref{thm:sgtr} is~\citep{fountoulakis2023graph},
which gives the qualitative guarantee for attention; we give a closed form for how
far a pruner of a given quality moves the threshold.

A parallel literature studies the limits of repeated aggregation. Over-smoothing
analyses~\citep{li2018deeper,oono2020graph} show node representations converge to
an uninformative subspace as depth grows, and spectral
accounts~\citep{nt2019revisiting} characterise graph convolution as a low-pass
filter. These results concern degradation with \emph{depth} under a fixed graph.
The effect we analyse is orthogonal: a single layer already suffers signal loss
when the \emph{topology itself} is adversarial. The two mechanisms compound, but
our critical ratio is present at $L = 1$.

\subsection{Heterophily in graph learning}

The finding that homophily-assuming GNNs degrade on heterophilic graphs motivated a
family of architectures that separate ego from neighbour representations, aggregate
over higher-order neighbourhoods, or combine low- and high-frequency
signals~\citep{zhu2020beyond,luan2022revisiting}. Ma et al.~\citeyearpar{ma2022homophily}
show in addition that homophily is not strictly necessary: GNNs can succeed on
heterophilic graphs when inter-class neighbourhood \emph{distributions} are
sufficiently distinguishable.

That observation matters for our setting, and we take a position on it. Under
adversarial camouflage the distinguishability Ma et al.\ rely on is precisely what
the adversary is optimising away: a bot that matches the human neighbourhood
distribution eliminates the signal their argument requires. The distinction between
benign heterophily---structured, and therefore exploitable---and adversarial
heterophily, engineered to be unexploitable, is in our view the conceptual point
separating bot detection from general heterophilic node classification. It is a
distinction about the \emph{provenance} of the inter-class edges, not about the
algebra of aggregation, which is why our threshold coincides with theirs in the
balanced case while the interpretation of the parameter differs.

\subsection{Robust aggregation and adversarial topology}
\label{sec:related-robust}

Camouflage is a contamination phenomenon, and two mature literatures have already
treated it as one.

The first replaces the mean with a robust estimator. In Byzantine-robust distributed
learning a fraction of workers report corrupted updates, and aggregation rules such
as Krum~\citep{blanchard2017machine} and coordinate-wise median or trimmed
mean~\citep{yin2018byzantine} bound the resulting error. The idea has been imported
into GNN neighbourhood aggregation directly: Geisler
et al.~\citeyearpar{geisler2020reliable} construct an aggregation with the maximal
breakdown point of $0.5$, so the bias stays bounded while fewer than half of a
node's edges are adversarial, and Chen et al.~\citeyearpar{chen2021understanding} attribute
the structural vulnerability of GCNs specifically to the non-robustness of the
weighted mean, analysed through the breakdown point.

We therefore avoid the term for our own threshold. A breakdown point asks when an
adversary can drive an estimate arbitrarily far under unbounded contamination;
$\rho^\star$ is instead a bias--variance crossover under \emph{bounded}
contamination---human features are not arbitrary---and marks where the $\sqrt{D}$
variance reduction ceases to pay for the $(1+c)\rho$ bias. The two quantities answer
different questions, and $\rho^\star$ is the weaker of the two.

The second literature attacks the topology. Nettack~\citep{zugner2018adversarial} and
Metattack~\citep{zugner2019adversarial} insert or remove edges to flip predictions,
and defences such as GNNGuard~\citep{zhang2020gnnguard} and
Pro-GNN~\citep{jin2020graph} filter them, largely heuristically. Two results in this
line bear directly on ours. Zhu et al.~\citeyearpar{zhu2022how} prove the bridge our model
assumes---that impactful structural attacks on homophilous data necessarily reduce
homophily---and derive the ego/neighbour-separation design principle from it.
Mujkanovic et al.~\citeyearpar{mujkanovic2022are} show empirically that under adaptive
attack, defended GNNs can underperform an MLP that ignores the topology outright;
our Corollary~\ref{cor:critical} can be read as a closed-form account of when that
should be expected. Certified-robustness work~\citep{bojchevski2019certifiable} is
orthogonal: it certifies that a prediction is unchanged within a perturbation
budget, never that the graph has become net-harmful beyond one. Gosch
et al.~\citeyearpar{gosch2023revisiting} come closest in combining a CSBM, an adversary,
and a Bayes-optimal yardstick, but use it to show GNNs are \emph{over}-robust,
insensitive past the point of genuine semantic change.

\subsection{Camouflage-resistant detectors}

CARE-GNN~\citep{dou2020caregnn} introduced the camouflage framing to GNN-based fraud
detection, using a label-aware similarity measure and a reinforcement-learning
neighbour selector to filter uninformative edges. Bot-detection work has since
adopted the pattern: BotHP~\citep{he2025bothp} separates ego and neighbour encoders
to mitigate interaction camouflage, and XHBot~\citep{dang2026xhbot} combines
spectral edge pruning, tri-channel aggregation, and contrastive disentanglement.

These methods are validated by ablation. An ablation establishes that removing a
component costs accuracy on a particular benchmark; it does not establish the
\emph{range of conditions} under which the component helps, nor whether a
multi-channel design is more expressive or merely better parameterised. The present
paper supplies exactly those missing statements for the three mechanisms XHBot
employs. We analyse XHBot's components because they instantiate the dominant design
pattern, not because the analysis is specific to that system: the edge-scoring
result applies to any feature-discrepancy pruning rule, and the expressivity result
to any aggregator that adds a permutation-invariant non-mean channel.

\end{document}
